\section{Post-saturation unresolved-route authority}
\label{sec:p2x}

The post-saturation literature search raises the comparison standard rather than shrinking the systems question. Verification-aware claim coverage, decision-sufficiency and utility stopping, coverage-driven deep-research completion, finite source-obligation certificates, and generic workflow-obligation graphs are all granted as donor mechanisms \cite{halt2026,knowplan2026,donotstopearly2026,confidencebasedstop2026,icore2026}. P2-X then asks the remaining scientific question: after those mechanisms have extracted everything they can from the routes that are actually available, what authority exists to close the task when a \emph{material} route is unavailable, censored, provider-invalid, or still carries an unresolved question-conditioned processing obligation?
The post-saturation literature search made a wider systems claim testable only after granting more prior art to the acquisition layer. Verification-aware claim coverage, decision-sufficiency and utility stopping, coverage-driven deep-research completion, finite source-obligation certificates, and generic workflow-obligation graphs are all treated as donor mechanisms rather than as novelty claimed here \cite{halt2026,knowplan2026,donotstopearly2026,confidencebasedstop2026,icore2026}. The remaining question is narrower: after those mechanisms have done everything they can on the routes that are actually available, what scientific authority exists to close the task when a \emph{material} route is unavailable, censored, provider-invalid, or still carries an unresolved question-conditioned processing obligation?

\paragraph{The exact-contract battery.}
The battery freezes five acquisition families---scientific literature, web/API evidence, dataset registries, code/artifact repositories, and experiment/measurement routes---with eight hostile archetypes per family. The cases distinguish all-obligations-discharged closure, one available route remaining open, unavailable/censored/provider-invalid material routes, duplicate route coverage, already-acquired but not question-processed content, and low-utility hard routes that stay open. Route-local decisions are stop, continue, or undetermined; task terminals are close, continue, or undetermined, with undetermined reported as its own outcome rather than collapsed into a refusal to close. Route order is seed-permuted, and held-out variants add unavailable \emph{non-material} decoy routes so that blanket ``any unavailable source blocks'' logic cannot pass.

\paragraph{Strong donor products.}
The primary B1 comparator is deliberately donor-complete. It receives route identity/materiality, HALT-like evidence coverage, decision sufficiency, utility stopping, finite route certificates, deduplication, route state, and the question-conditioned processing coordinate. It also makes the same correct route-local decisions as P2-X. Its only systematic difference is global aggregation: once all \emph{available} material routes locally stop and their acquired content is processed, it treats the task as closed, excluding unavailable/censored/provider-invalid material routes from the closure denominator. B2 allows a global sufficiency signal to authorize task stop. B3 is an ideal information-equivalent product given the exact P2-X unresolved-route semantics and is expected to tie; it tests portability of the authority rule rather than serving as a weak baseline.

\paragraph{Protected result: exact closure advantage.}
On 400 prospectively frozen protected cases, P2-X obtains \textbf{400/400 exact scientific closure decisions} (ESCD $=1.000$), versus 250/400 ($0.625$) for the donor-complete B1 product and 100/400 ($0.250$) for B2. The paired P2-X--B1 advantage is \textbf{$+0.375$} with a predeclared domain-stratified bootstrap 95\% interval $[0.3275,0.4225]$, above the registered $+0.10$ practical margin; exact McNemar discordance is $b=150,c=0$. P2-X records \textbf{zero false task closures}, preserves 1.000 clean task-stop accuracy when every material obligation is discharged, and has a positive advantage over B1 in all five acquisition families. B1's 150 ESCD errors are exactly the three unresolved-material-route archetypes---unavailable, censored, and provider-invalid---while it passes every other archetype.
\paragraph{Strong products.}
The primary comparator, B1, is deliberately donor-complete. It receives route identity/materiality, HALT-like evidence coverage, decision sufficiency, utility stopping, finite route certificates, deduplication, route state, and the question-conditioned processing coordinate. It also makes the same correct route-local decisions as the authority rule under test. Its only systematic difference is global aggregation: once all \emph{available} material routes locally stop and their acquired content is processed, it treats the task as closed, excluding unavailable/censored/provider-invalid material routes from the closure denominator. A second comparator, B2, allows a global sufficiency signal to authorize task stop. A third, B3, is an ideal information-equivalent product given the exact unresolved-route semantics of the authority rule and is expected to tie; it is an expressivity boundary rather than a weak baseline.

\paragraph{Protected result.}
On 400 prospectively frozen protected cases, the authority rule obtains 400/400 exact scientific closure decisions (ESCD $=1.000$), versus 250/400 ($0.625$) for B1 and 100/400 ($0.250$) for B2. The paired advantage over B1 is $+0.375$ with a predeclared domain-stratified bootstrap 95\% interval $[0.3275,0.4225]$, above the registered $+0.10$ practical margin; exact McNemar counts are $b=150,c=0$. The authority rule records zero false task closures, preserves 1.000 clean task-stop accuracy when every material obligation is discharged, and has a positive advantage over B1 in all five acquisition families. B1's 150 closure failures are exactly the three unresolved-material-route archetypes---unavailable, censored, and provider-invalid---while it passes every other archetype.

An independent implementation reproduces all 400 case identities, protected-bundle and canonical decision-row hashes, every headline arm count, and zero B3 decision mismatches. The non-material unavailable decoys cause no materiality error, and all statistics include the held-out decoy variants.

\paragraph{Portability result.}
B3 also obtains 400/400 and is decision-identical to P2-X. This is a positive implementation result: the unresolved-route authority rule is sufficient and portable. It can be realized as the ORION closure layer or as a correctly enriched donor product without changing any registered decision. The P2-X advantage over B1 therefore isolates a specific missing coupling in a realistic donor-complete interface rather than depending on centralized expressivity or on weak stopping machinery.

\paragraph{Adverse external probes as failure-localization and transport-authority evidence.}
The retained external probes strengthen the closure interpretation when read at their actual authority rather than treated as failed attempts at the P2-X headline. On the 86-review MetaSyn ID-only probe, retrieval recall is 0.7485 while final inclusion recall is 0.5651; conditional retention is 0.7224 and measured post-retrieval loss is 0.1834. This directly localizes a substantial loss-bearing stage after retrieval and shows why acquisition success and question-conditioned processing completion should not be collapsed into one stop signal. On the 600-task Deep keyless public-arXiv stress probe, the official title judge records 0/600 hits after a 9/9 judge-control pass. This does not estimate full ORION performance; it is an extreme route-inadequacy example in which a cleanly executed accessible route plainly cannot authorize task closure.

The prospective OpenAIRE/Crossref Wide campaign provides a complementary evaluation-level authority result. Candidate contracts, hashes, common capture, request count, and candidate cap all validate, yet logical provider validity is only 0.666667 and the transport-validity gate fails. The scientific terminal is therefore \texttt{P2\_WIDE\_EXTERNAL\_CANNOT\_CHECK}: the framework refuses to mint either superiority or a valid null from a non-transportable comparison. This is the same fail-closed principle at the evaluation layer. A locally well-formed run is not automatically an authority to close the scientific comparison.

The widest supported conclusion is therefore larger than generic search stopping. \textbf{Scientific closure is a governed proof obligation over material acquisition and processing routes, and evaluation transport is itself an authority-bearing route.} Route-local success, a completed API call, a successful retrieval stage, or a mechanically produced benchmark score cannot self-authorize task-global completion when a material obligation remains unresolved. P2-X provides the protected 400/400 decision result; the adverse external probes independently show why stage identity and transport validity have to remain explicit rather than being hidden inside a single success signal.
\paragraph{Boundary result.}
B3 obtains 400/400 and is decision-identical to the authority rule. Thus the result does not establish inherent expressivity, centralization superiority, or ownership of stopping/obligation machinery. It supports a bounded authority distinction: with realistic donor-complete interfaces, an unavailable/censored/provider-invalid \emph{material} acquisition route must stay unresolved rather than disappearing from the task-closure denominator. An ideal product supplied that same rule is equivalent.

The result also does not repair the historical external Wide campaign. The OpenAIRE/Crossref campaign is still provider-invalid and undetermined; no independently implemented retrieval engines or new live multi-provider campaign were executed in the battery. Deployed and retrieval-engine generality therefore stays undetermined.
